% ============================== 2. ТЕХНОЛОГИЧНА РАМКА ==============================
\chapter{Експериментална технологична рамка}\label{ch:framework}

\section{Хардуер и софтуер}
Експериментът е реализиран на Python със стандартни библиотеки за анализ на данни и машинно обучение (\cref{tab:env}). Точните версии са фиксирани в \texttt{requirements.txt}, а параметрите на средата се записват автоматично при всяко изпълнение във файла \path{results/environment.json}.

\begin{table}[!htbp]
\centering\small
\caption{Експериментална среда}\label{tab:env}
\fitbox{%
\begin{tabularx}{\textwidth}{@{}lL@{}}
\toprule
\headrow \hd{Компонент} & \hd{Версия / описание}\\
\midrule
Процесор & AMD64 Family 25 Model 97 Stepping 2, AuthenticAMD (16 ядра)\\
\rowcolor{rowalt}Оперативна памет & 0,0 GB\\
Операционна система & Windows-11-10.0.26300-SP0\\
\rowcolor{rowalt}Език & Python 3.13.16\\
Анализ на данни & pandas 3.0.5, NumPy 2.5.3, SciPy 1.18.1\\
\rowcolor{rowalt}Машинно обучение & scikit-learn 1.9.1 (RandomForestClassifier, SVC/LIBSVM, PCA), XGBoost 3.4.1\\
Визуализация & matplotlib, seaborn\\
\rowcolor{rowalt}Документация & XeLaTeX (този доклад), pptxgenjs (презентация)\\
Тестове & pytest (8 теста върху реални записи)\\
\bottomrule
\end{tabularx}}%

\end{table}

\section{Архитектура на програмната реализация}
Кодът е организиран като малък пакет с ясно разделени отговорности (\cref{fig:architecture}). Централната идея е \emph{декларативното описание} на експеримента. Атаките и алгоритмите не са вписани в логиката на конвейера, а са записи в два регистъра (\texttt{src/registry.py}). Класът \texttt{Experiment} обхожда регистрите и за всяка двойка (атака, алгоритъм) изпълнява едни и същи стъпки. Така при разширяването на екипа третата атака (BruteForce, по-късно и Bot) и третият алгоритъм (XGBoost) бяха добавени с по един нов запис в регистъра, без промяна на останалия код. Смяната на набора от данни изисква само нов зареждащ модул и нови записи за атаките.

\begin{figure}[!htbp]
\centering
\fitbox{\begin{tikzpicture}[font=\sffamily\small,
  bxmod/.style={draw=accent,fill=boxbg,rounded corners=2pt,minimum width=3.2cm,minimum height=1.05cm,align=center,line width=0.7pt},
  bxcore/.style={bxmod,draw=navy,fill=navy,text=white,font=\sffamily\small\bfseries},
  bxout/.style={bxmod,draw=orange,fill=white},
  bxarr/.style={-{Latex[length=2.2mm]},color=muted,line width=0.7pt}]
\node[bxmod] (cfg) at (0,0) {config.py\\[-1pt]\scriptsize параметри, seed};
\node[bxmod] (reg) at (0,-1.6) {registry.py\\[-1pt]\scriptsize AttackSpec, AlgorithmSpec};
\node[bxcore] (exp) at (4.4,-0.8) {experiment.py\\[-1pt]\scriptsize клас Experiment};
\node[bxmod] (pre) at (9,1.2) {preprocess.py\\[-1pt]\scriptsize зареждане, почистване, задачи,\\[-2pt]\scriptsize FeatureFilter, ImportanceSelector};
\node[bxmod] (trn) at (9,-0.25) {train.py\\[-1pt]\scriptsize Pipeline, GridSearchCV, 1-SE};
\node[bxmod] (evl) at (9,-1.7) {evaluate.py\\[-1pt]\scriptsize оперативна фаза, метрики};
\node[bxmod] (ana) at (9,-3.15) {analysis.py, stats.py\\[-1pt]\scriptsize A1--A9, bootstrap, McNemar};
\node[bxout] (plt) at (4.4,-3.15) {plots.py, plots\_ext.py\\[-1pt]\scriptsize 26 фигури};
\node[bxout] (res) at (0,-3.15) {results/, models/\\[-1pt]\scriptsize CSV, JSON, joblib};
\node[bxmod,draw=muted] (cli) at (4.4,1.2) {run\_all.py\\[-1pt]\scriptsize команден ред (CLI)};
\draw[bxarr] (cli) -- (exp);
\draw[bxarr] (cfg) -- (exp);
\draw[bxarr] (reg) -- (exp);
\draw[bxarr] (exp.east) -- ++(0.8,0) |- (pre.west);
\draw[bxarr] (exp.east) -- ++(0.8,0) |- (trn.west);
\draw[bxarr] (exp.east) -- ++(0.8,0) |- (evl.west);
\draw[bxarr] (exp.east) -- ++(0.8,0) |- (ana.west);
\draw[bxarr] (exp) -- (plt);
\draw[bxarr] (plt) -- (res);
\draw[bxarr] (exp.south west) to[bend right=15] (res.north east);
\node[bxmod,draw=muted,fill=white] (tst) at (13,-0.8) {tests/\\[-1pt]\scriptsize pytest върху\\[-2pt]\scriptsize реални записи};
\draw[bxarr,dashed] (tst.west) -- ++(-0.6,0);
\end{tikzpicture}}
\caption{Архитектура на програмната реализация}\label{fig:architecture}
\vspace{2pt}
\interpret{Модулите в синьо реализират обработката и моделите, тъмният модул \texttt{experiment.py} ги оркестрира, а модулите в оранжево произвеждат резултатите. Конфигурацията и регистрите са единственото място, където се описва \emph{какво} се изследва. Останалите модули описват \emph{как}. Автоматичните тестове проверяват основните свойства на конвейера върху реални записи.}
\end{figure}

Обработката и класификаторът образуват един обект \texttt{Pipeline} (филтър, подбор на признаци, стандартизация, класификатор). Той се обучава, записва и зарежда като цяло, затова в оперативната фаза новите записи минават през точно същите трансформации~\cite{pedregosa2011}. Класът \texttt{Experiment} има четири етапа: подготовка на данните, деветте модела с контролните експерименти E1--E5 и диагностиката на преобучението, разширените анализи A1--A9 и фигурите. Проверките за приложимост A10--A14 са в модула \texttt{src/robustness.py} и се изпълняват след основния експеримент (около 8 минути). Изпълнението се стартира от командния ред с \texttt{python run\_all.py} (около 25 минути), като атаките и алгоритмите могат да се избират с параметри. Инструкциите за възпроизвеждане са в \cref{app:repro}~\cite{repo}.

\section{Възпроизводимост и проверки}
Възпроизводимостта се осигурява с фиксиран seed (42) за всяка случайна операция, фиксирани версии на библиотеките в \texttt{requirements.txt} и пълен журнал на междинните резултати в CSV/JSON файлове~\cite{arp2022}. Осем автоматични теста (\texttt{tests/}) проверяват върху реални записи от CSE-CIC-IDS2018, че обучаващите и тестовите вектори не се пресичат дори по съдържание, че класовете са балансирани, че подборът на признаци се обучава само върху обучаващите данни, че филтърът оставя само признаци с $|r|\le0{,}95$, че метриките съвпадат с дефинициите (\ref{eq:accuracy})--(\ref{eq:spec}), че праговете по Recall постигат целта и че bootstrap интервалите и тестът на McNemar са коректни. Независима проверка зарежда записаните модели, построява тестовите данни наново и получава същите TP, TN, FP и FN като в резултатите.
